\documentclass[12pt]{book}
\usepackage[utf8]{inputenc}
\usepackage[margin=1in]{geometry}
\usepackage{amsmath}
\usepackage{amssymb}

\title{Quantum Theories: The Missing Big Picture}
\author{A Unified Framework}
\date{\today}

\begin{document}

\maketitle

\section*{Blurb}

Quantum theory books abound, but readers are left in a fragmented maze. This course provides a unified framework by emphasizing quantum theories' goal: predict time evolution of states and observables.

\section*{The 11-Dimensional Framework}

Time evolution equations map onto an 11-dimensional space:

1. Formalism (Operator vs Path Integral)
2. Relativity (Non-relativistic, Special, General)
3. Particle Number (Single, Fixed-N, Fock)
4. Picture (Schrödinger, Heisenberg, Interaction)
5. State Purity (Pure, Mixed)
6. System (Closed, Open)
7. Spin (Spinless, 1/2, Arbitrary)
8. Mass (Massive, Massless)
9. Approximation (Exact, Perturbative, Linearized)
10. Classical Limit (Quantum, Semiclassical, Classical)
11. Solutions (Analytical, Approximate, Numerical)

\tableofcontents

\part{Foundations}

\chapter{The Universal Language}

\section{Why Time Evolution?}

The central question of quantum mechanics is: How does a quantum system change with time?

Core principle: All quantum theories answer this through time evolution equations:
$$i\hbar\frac{d}{dt}|\psi(t)\rangle = \hat{H}|\psi(t)\rangle$$

Variations in this equation form the 11-dimensional landscape.

\section{The 11 Dimensions}

\subsection{Dimension 1: Formalism}

How we express time evolution:
\begin{itemize}
  \item Operator formalism (differential equations)
  \item Path integral formalism (functional integrals)
\end{itemize}

\subsection{Dimension 2: Relativity}

Spacetime symmetries:
\begin{itemize}
  \item Non-relativistic (Galilean)
  \item Special relativistic (Lorentz)
  \item General relativistic (Diffeomorphism)
\end{itemize}

\subsection{Dimension 3: Particle Number}

Number of particles:
\begin{itemize}
  \item Single particle
  \item Fixed-N system
  \item Fock space (variable number)
\end{itemize}

\subsection{Dimension 4: Picture}

Where time dependence lives:
\begin{itemize}
  \item Schrödinger (states evolve)
  \item Heisenberg (operators evolve)
  \item Interaction (both evolve)
\end{itemize}

\subsection{Dimension 5: State Purity}

Knowledge completeness:
\begin{itemize}
  \item Pure states
  \item Mixed states
\end{itemize}

\subsection{Dimension 6: System}

Boundary conditions:
\begin{itemize}
  \item Closed (isolated)
  \item Open (environmental coupling)
\end{itemize}

\subsection{Dimension 7: Spin}

Internal angular momentum:
\begin{itemize}
  \item Spinless
  \item Spin-1/2
  \item Arbitrary spin
\end{itemize}

\subsection{Dimension 8: Mass}

Particle properties:
\begin{itemize}
  \item Massive
  \item Massless
\end{itemize}

\subsection{Dimension 9: Approximation}

Solution strategy:
\begin{itemize}
  \item Exact
  \item Perturbative
  \item Linearized
\end{itemize}

\subsection{Dimension 10: Classical Limit}

Quantum-classical transition:
\begin{itemize}
  \item Quantum
  \item Semiclassical
  \item Classical
\end{itemize}

\subsection{Dimension 11: Solutions}

Implementation:
\begin{itemize}
  \item Analytical
  \item Approximate
  \item Numerical
\end{itemize}

\chapter{The Simplest Case}

\section{Non-Relativistic Spinless Schrödinger Equation}

The Elementary Schrödinger Equation:
$$i\hbar\frac{\partial\psi(\mathbf{r},t)}{\partial t} = -\frac{\hbar^2}{2m}\nabla^2\psi(\mathbf{r},t) + V(\mathbf{r},t)\psi(\mathbf{r},t)$$

Dimensions: Non-relativistic, Single particle, Schrödinger picture, Pure state, Closed system, Spinless, Massive.

\subsection{Operator Formalism}

Abstract form:
$$i\hbar\frac{d|\psi(t)\rangle}{dt} = \hat{H}|\psi(t)\rangle$$

where the Hamiltonian is:
$$\hat{H} = \frac{\hat{p}^2}{2m} + V(\hat{\mathbf{r}},t)$$

\subsection{Formal Solution}

For time-independent Hamiltonians:
$$|\psi(t)\rangle = e^{-i\hat{H}t/\hbar}|\psi(0)\rangle$$

The time evolution operator $U(t) = e^{-i\hat{H}t/\hbar}$ is unitary and preserves probability.

\subsection{Energy Eigenstates}

When $V$ is time-independent, expand in energy eigenstates:
$$\hat{H}|n\rangle = E_n|n\rangle$$

yielding:
$$|\psi(t)\rangle = \sum_n c_n e^{-iE_n t/\hbar}|n\rangle$$

where $c_n = \langle n|\psi(0)\rangle$.

\section{The Hydrogen Atom}

\subsection{Hamiltonian}

For an electron orbiting a proton:
$$\hat{H} = \frac{\hat{p}^2}{2m_e} - \frac{ke^2}{r}$$

\subsection{Energy Spectrum}

The exact solution:
$$E_n = -\frac{13.6 \text{ eV}}{n^2}, \quad n = 1,2,3,\ldots$$

\chapter{Mathematical Foundations}

\section{Hilbert Space}

Quantum states live in a complex Hilbert space with:
\begin{itemize}
  \item Inner product: $\langle\phi|\psi\rangle$
  \item Linearity
  \item Completeness: $\sum_n |n\rangle\langle n| = I$
\end{itemize}

\section{Operators}

\subsection{Hermitian Operators}

Observables: $\hat{A}^\dagger = \hat{A}$

Properties:
\begin{itemize}
  \item Real eigenvalues
  \item Orthogonal eigenstates
  \item Complete basis
\end{itemize}

\subsection{Unitary Operators}

Evolution: $\hat{U}^\dagger\hat{U} = I$

\section{Commutators and Uncertainty}

Canonical commutation relations:
$$[\hat{x}_i, \hat{p}_j] = i\hbar\delta_{ij}$$

Heisenberg uncertainty principle:
$$\Delta x \cdot \Delta p \geq \frac{\hbar}{2}$$

\part{Dimension Variations}

\chapter{Pictures}

\section{Heisenberg Picture}

States static, operators evolve:
$$|\psi_H(t)\rangle = |\psi_H(0)\rangle$$
$$\hat{A}_H(t) = \hat{U}^\dagger(t)\hat{A}_S\hat{U}(t)$$

\subsection{Equation of Motion}

$$\frac{d\hat{A}_H}{dt} = \frac{i}{\hbar}[\hat{H}, \hat{A}_H]$$

\section{Interaction Picture}

For perturbations: $\hat{H} = \hat{H}_0 + \hat{H}_I(t)$

$$|\psi_I(t)\rangle = \hat{U}_0^\dagger(t)|\psi_S(t)\rangle$$

Starting point for perturbation theory.

\chapter{Relativistic Quantum Mechanics}

\section{Klein-Gordon Equation}

From $E^2 = \mathbf{p}^2c^2 + m^2c^4$:
$$\left(\frac{1}{c^2}\frac{\partial^2}{\partial t^2} - \nabla^2 + \frac{m^2c^2}{\hbar^2}\right)\phi = 0$$

Describes spin-0 particles.

\section{The Dirac Equation}

Incorporating spin-1/2:
$$i\hbar\frac{\partial\psi}{\partial t} = (\boldsymbol{\alpha} \cdot \mathbf{p}c + \beta mc^2)\psi$$

Naturally incorporates antiparticles.

\chapter{Quantum Field Theory}

\section{From Particles to Fields}

Second quantization:
$$\hat{\phi}(\mathbf{r}) = \int \frac{d^3p}{(2\pi)^{3/2}\sqrt{2E_\mathbf{p}}}(\hat{a}_\mathbf{p}e^{i\mathbf{p} \cdot \mathbf{r}} + \hat{a}_\mathbf{p}^\dagger e^{-i\mathbf{p} \cdot \mathbf{r}})$$

\section{Lagrangian Formalism}

The action:
$$S = \int d^4x \, \mathcal{L}(\phi, \partial_\mu\phi)$$

Canonical quantization:
$$[\hat{\phi}(\mathbf{x}), \hat{\pi}(\mathbf{x}')] = i\hbar\delta^3(\mathbf{x} - \mathbf{x}')$$

\chapter{Quantum Electrodynamics}

\section{Electron-Photon Coupling}

$$\mathcal{L} = \bar{\psi}(i\gamma^\mu(\partial_\mu - ieA_\mu) - m)\psi - \frac{1}{4}F_{\mu\nu}F^{\mu\nu}$$

Coupling constant: $\alpha \approx 1/137$

\section{QED Precision}

\begin{itemize}
  \item Lamb shift: 1.06 GHz (theory) vs 1.057 GHz (experiment)
  \item Electron g-factor: agreement to 1 part in $10^{12}$
\end{itemize}

\chapter{Open Systems}

\section{Density Matrix}

For mixed states:
$$\rho = \sum_i p_i |\psi_i\rangle\langle\psi_i|$$

Expectation value:
$$\langle A \rangle = \text{Tr}[\rho\hat{A}]$$

\section{Master Equations}

Open system evolution:
$$\frac{d\rho_S}{dt} = -\frac{i}{\hbar}[\hat{H}_S, \rho_S] + \sum_\mu \left(\hat{L}_\mu\rho_S\hat{L}_\mu^\dagger - \frac{1}{2}\{\hat{L}_\mu^\dagger\hat{L}_\mu, \rho_S\}\right)$$

\part{Solution Techniques}

\chapter{Perturbation Theory}

\section{Time-Independent}

When $\hat{H} = \hat{H}_0 + \lambda\hat{H}_1$:

$$E_n = E_n^{(0)} + \lambda E_n^{(1)} + \lambda^2 E_n^{(2)} + \cdots$$

\section{First-Order Corrections}

$$E_n^{(1)} = \langle n^{(0)}|\hat{H}_1|n^{(0)}\rangle$$

$$|n^{(1)}\rangle = \sum_{k \ne n} \frac{\langle k^{(0)}|\hat{H}_1|n^{(0)}\rangle}{E_n^{(0)} - E_k^{(0)}}|k^{(0)}\rangle$$

\chapter{Path Integrals}

\section{Feynman Approach}

$$\langle\mathbf{x}_f, t_f | \mathbf{x}_i, t_i\rangle = \int \mathcal{D}\mathbf{x} \, e^{iS[\mathbf{x}]/\hbar}$$

where $S[\mathbf{x}] = \int dt \left[\frac{m}{2}\dot{\mathbf{x}}^2 - V(\mathbf{x})\right]$ is the action.

\section{WKB Approximation}

$$\psi(x) = A(x)e^{iS(x)/\hbar}$$

\chapter{Numerical Methods}

\section{Imaginary Time Evolution}

To find ground state, evolve in imaginary time $\tau = it$:
$$\frac{\partial}{\partial\tau}\psi = -\hat{H}\psi$$

\section{Quantum Computing}

Qubits: $|\psi\rangle = \alpha|0\rangle + \beta|1\rangle$

Key algorithms:
\begin{itemize}
  \item Grover: Search in $\sim\sqrt{N}$ queries
  \item Shor: Factor in polynomial time
\end{itemize}

\part{Synthesis}

\chapter{Connecting the Dimensions}

\section{Taxonomy of Quantum Theories}

\begin{tabular}{|l|l|}
\hline
\textbf{Theory} & \textbf{Dimensions} \\
\hline
Quantum Mechanics & Non-rel, Single, Pure \\
\hline
QFT & Relativistic, Fock \\
\hline
QED & Special rel, Fock, Massive \\
\hline
QCD & Special rel, Fock, Massless \\
\hline
Open Systems & Open, Mixed \\
\hline
Quantum Info & Computational \\
\hline
\end{tabular}

\section{The Schrödinger Equation as Lodestar}

All theories ultimately answer: how do states and observables evolve?

$$i\hbar\frac{d}{dt}|\psi(t)\rangle = \hat{H}(t)|\psi(t)\rangle$$

Variations:
\begin{itemize}
  \item Different Hamiltonians
  \item Different Hilbert spaces
  \item Different representations
  \item Different solution methods
\end{itemize}

\section{Why These 11 Dimensions?}

D1-D2-D3: Physical content

D4-D5-D6: Formalism

D7-D8: Particle properties

D9-D10-D11: Solution strategy

\chapter{Open Problems}

\section{Quantum Gravity}

Quantize gravity without infinities.

Approaches:
\begin{itemize}
  \item String Theory
  \item Loop Quantum Gravity
  \item Causal Sets
\end{itemize}

\section{Extreme Conditions}

Quark-Gluon Plasma, Black Hole Information, Holography.

\chapter{Conclusion}

\section{From Fragmentation to Unity}

Students see quantum mechanics as disconnected topics. The unifying principle is time evolution of states and observables.

\section{The 11-Dimensional Framework}

Provides:
\begin{enumerate}
  \item Conceptual scaffolding
  \item Predictive power
  \item Problem-solving tools
  \item Research guidance
\end{enumerate}

\section{Final Thought}

Quantum mechanics is not mysterious. The "missing big picture" is that quantum theories form a coherent whole once viewed through time evolution in an 11-dimensional space.

The Schrödinger equation is the starting point from which all quantum theories emerge through systematic variation.

\vspace{2cm}
\textit{The quantum world awaits those with the right framework to see it clearly.}

\appendix

\chapter{Key Equations}

Schrödinger: $i\hbar\frac{\partial\psi}{\partial t} = \hat{H}\psi$

Commutation: $[\hat{x}_i, \hat{p}_j] = i\hbar\delta_{ij}$

Uncertainty: $\Delta x \cdot \Delta p \geq \hbar/2$

Klein-Gordon: $(\Box + m^2)\phi = 0$

Dirac: $(i\gamma^\mu\partial_\mu - m)\psi = 0$

Partition function: $Z = \text{Tr}[e^{-\beta\hat{H}}]$

\end{document}
